% vim: tw=78 ai
\documentclass{article}

%\usepackage{a4wide}
\usepackage[utf8x]{inputenc}
\usepackage[czech]{babel}
\usepackage[IL2]{fontenc}

\pagestyle{empty}

\begin{document}
\begin{flushright}
\today
\end{flushright}
\begin{center}
\Large\textbf{Posudek vedoucího na bakalářskou práci \\
\normalsize%
\textit{Tomáš Jamriško:} Vizualizácia modelovania terénu}
\end{center}

Dle zadání měla bakalářská práce nastudovat problematiku počítačového
modelování terénu a implementovat SW generující 3D model z mapových podkladů.

Text práce je na 26 stranách strukturovaný do 8 kapitol. V úvodu a druhé
kapitole se čtenář dozví motivaci práce. V kapitolách 3-6 jsou popsány
principy, na kterých má práce stavět. V 7. kapitole je uvedena samotná
implementace. Práci ukončuje závěr.

Struktura práce je sice dobrá, ale zasloužila by si lepší provázání kapitol,
aby bylo jasnější, k čemu v práci jsou. Taktéž chybí nějaké faktické porovnání
nastudovaných metod a komplexnější zdůvodnění, proč byly v implementaci
použity právě vybrané metody. I přesto je text kapitol 3-6 dostatečný k tomu,
aby splnil zadání.

7. kapitola se věnuje implementační části. Jsou zde uvedeny použité knihovny a
služby poskytující výškové podklady. Část textu této kapitoly je psána
nevhodně 1. osobou.

V celém textu práce měl být, dle požadavků na bakalářskou práci, kladen větší
důraz na přínos práce na úkor např. popisu známých technik. Autor měl klást
větší důraz kupříkladu na návrh aplikace, jak bylo uvedeno v zadaní.

Předložená implementační část ve formě spustitelného souboru pro Windows není
funkční. Nelze terén načíst ani zpracovat. Po překladu přiložených zdrojových
souborů na OS Linux aplikace dokáže modely načítat, zpracovat, zobrazovat a
ukládat, nicméně působí nepropracovaně, přičemž nejsou výjimkou pády celé
aplikace. Aplikace není příliš uživatelsky přívětivá a v textu práce chybí
jednoduchý návod, jak s ní pracovat.

Po jazykové stránce obsahuje práce jen několik překlepů a typografické chyby
-- až příliš částé přetoky řádků.

Většina z výše uvedených nedostatků by byla vymícena, byla-li by práce
konzultována s vedoucím těsně před jejím odevzdáním, k čemuž bohužel nedošlo.
Celkově práce \textbf{splňuje zadání i požadavky na bakalářskou práci}.
Navrhuji práci \textbf{uznat} jako bakalářskou a vzhledem k výše uvedeným
nedostatkům hodnotit stupněm C.

\begin{flushright}
Jiří Slabý
\end{flushright}

\end{document}
